
\documentclass[10pt]{article}
\usepackage[utf8]{inputenc}
\usepackage[a4paper, left=2.54cm, right=2.54cm, top=3cm, bottom=3cm]{geometry} % Márgenes personalizados a 2cm
\usepackage[spanish]{babel} % Para fechas y formato en español
\usepackage{multicol}
\usepackage{enumitem}
\usepackage{fancyhdr}
\usepackage{titlesec}
\usepackage{amsmath}

% Configuración de página
\pagestyle{fancy}
\fancyhf{} % Limpia encabezados y pies de página
\fancyhead[L]{CEBA - Santo Domingo}
\fancyhead[R]{\today}
\fancyfoot[C]{\thepage} % Número de página centrado en el pie
\renewcommand{\headrulewidth}{0.8pt}

% Configuración de títulos y espaciado
\titleformat{\section}[block]{\large\bfseries\centering}{\thesection}{1em}{}
\setlength{\columnsep}{1cm}
\setlength{\parskip}{0.8em}
\setlength{\parindent}{0em}


\begin{document}

	\begin{center}
		\Large\bfseries Productos Notables
	\end{center}

\begin{multicols}{2}
	\section*{Ejercicios}
		\begin{enumerate}[label=\textbf{\arabic*.}, itemsep=0.4cm]
			\item Calcula $a^2 + b^2$ si: \\
			$a + b = 5$, $a \cdot b = 3$.

			\item Calcula $ab$ si: \\
			$a - b = 5$, $a^2 + b^2 = 17$.

			\item Calcula $a - b$ si: \\
			$a + b = 4$, $a^2 + b^2 = 9$, $a < b$.

			\item Calcula $a - b$ si: \\
			$a + b = 6$, $a \cdot b = 6$, $a < b$.

			\item Calcula $a - b$ si: \\
			$a + b = 4$, $a \cdot b = 2$.

			\item Simplifica: \\
			$\dfrac{(x + 1 + 2\sqrt{x})(x + 1 - 2\sqrt{x}) + (x + 1)^2}{2(x + 1 + \sqrt{2x})(x + 1 - \sqrt{2x})}$.

			\item Dado $\dfrac{a}{b} + \dfrac{b}{a} = 1$ ($a, b \neq 0$), calcula: \\
			$\dfrac{a^4 + b^4}{a^2 \cdot b^2}$.

			\item Determina la expresión simplificada de: \\
			$(a^b + a^{-b})(a^b - a^{-b})(a^{4b} + 1 + a^{-4b})$.

			\item Reduce la siguiente expresión: \\
			$A = (x + 1)(x - 1)(x^2 - x + 1)(x^2 + x + 1) - x^6$.

			\item Determina el valor de $abc$ si: \\
			$a + \dfrac{1}{b} = 1$, $b + \dfrac{1}{c} = 1$, con $b \neq 0$, $c \neq 0$.

			\item Si $x - x^{-1} = 1$ ($x \neq 0$), entonces los valores de: \\
			$x^2 + x^{-2}$ y $x^3 - x^{-3}$.

			\item Determina el valor de $x$ que verifica: \\
			$\sqrt[3]{14 + \sqrt{x}} + \sqrt[3]{14 - \sqrt{x}} = 4$.

			\item Determina el valor de $x$ que verifica: \\
			$\sqrt{3 + \sqrt{x}} + \sqrt{3 - \sqrt{x}} = 2\sqrt{2}$.

			\item Sean los números: \\
			$a = \dfrac{1}{2} + \dfrac{1}{\sqrt{2}}$, $b = \dfrac{1}{2} - \dfrac{1}{\sqrt{2}}$. \\
			Entonces, calcula $a + \dfrac{1}{a} + b + \dfrac{1}{b}$.


			% Para añadir más ejercicios, simplemente copia y pega el siguiente bloque:
			% \item [Enunciado del ejercicio]
		\end{enumerate}

\end{multicols}

\end{document}